% =============================================================================
% apendiceC.tex — ANÁLISE ESTÁTICA DA VIGA PRINCIPAL
% GERADO AUTOMATICAMENTE por mlc_projeto/output/apendice_analise.py a cada
% execução do pipeline (issue 51). NÃO EDITAR À MÃO.
% =============================================================================

\chapter{ANÁLISE ESTÁTICA DA VIGA PRINCIPAL}\label{ap:analise-estatica}

Este apêndice reúne os resultados da análise estática da viga principal, biapoiada e sem empuxo, descrita no \autoref{cap:modelagem}, para todas as combinações de estado-limite último: 14 com vento dominante sob $1{,}35\,G + 1{,}5\,W$, 14 com vento de sucção sob $0{,}9\,G + 1{,}5\,W$, uma para cada padrão de coeficientes de pressão, e a de sobrecarga dominante $1{,}35\,G + 1{,}5\,Q$. Os esforços são referidos ao eixo da peça: o cortante $V$ é perpendicular ao eixo e o esforço normal $N$ é positivo em compressão. Os valores são os da execução do programa que produziu o memorial, de onde este apêndice é gerado.

\section{COMBINAÇÃO GOVERNANTE}\label{sec:an-governante}

A combinação 29, $1{,}35\,G + 1{,}5\,Q$, governa a flexão, o cortante e o esforço normal da viga, com momento máximo de 505~kN$\cdot$m a 10,2~m do apoio. A \autoref{fig:an-mvn-governante} apresenta os seus diagramas.

\figura
{Diagramas de esforços da combinação governante}
{0.65}
{fig/analise_MVN_governante.pdf}
{Elaborado pelo autor (2026)}
{an-mvn-governante}
{}
{}

\section{TODAS AS COMBINAÇÕES}\label{sec:an-todas}

A \autoref{fig:an-mvn-todas} sobrepõe os diagramas das 29 combinações, e a \autoref{tab:an-esforcos} lista os esforços de cada uma. O momento varia entre $-$391~kN$\cdot$m, na combinação 17, e 505~kN$\cdot$m, na combinação 29. A combinação 15 é a de maior sucção e governa o arrancamento das bases e a flambagem lateral com o intradorso comprimido.

\figura
{Diagramas de esforços de todas as combinações de estado-limite último}
{0.85}
{fig/analise_MVN_combinacoes.pdf}
{Elaborado pelo autor (2026)}
{an-mvn-todas}
{}
{}

\tabela{Esforços na viga principal por combinação de estado-limite último}{\scriptsize\setlength{\tabcolsep}{2.6pt}
\begin{tabular}{@{}rllrrrrrrrrr@{}}
\toprule
    \textbf{Nº} & \textbf{Padrão de vento} & \textbf{$c_{pi}$} & \textbf{$q_{esq}$} & \textbf{$q_{dir}$} & \textbf{$M_{max}$} & \textbf{$x$} & \textbf{$M_{ap,d}$} & \textbf{$R_A$} & \textbf{$R_B$} & \textbf{$V_{max}$} & \textbf{$N_{max}$} \\
    & & & \multicolumn{2}{c}{(kN/m)} & (kN$\cdot$m) & (m) & (kN$\cdot$m) & (kN) & (kN) & (kN) & (kN) \\
\midrule
    \multicolumn{12}{@{}l}{\textit{Vento dominante, $1{,}35\,G + 1{,}5\,W$}} \\
    1 & $\alpha_v = 0^\circ$, t1, esq. & $+0{,}8$ & $-$7,73 & $-$4,14 & $-$281 & 8,71 & $-$275 & $-$77,5 & $-$50,8 & $-$69,1 & $-$6,09 \\
    2 & $\alpha_v = 0^\circ$, t1, esq. & $-0{,}3$ & 0,179 & 3,78 & 136 & 11,9 & 128 & 11,9 & 29,1 & $-$28,9 & 2,55 \\
    3 & $\alpha_v = 0^\circ$, t1, dir. & $+0{,}8$ & $-$4,14 & $-$7,73 & $-$283 & 11,5 & $-$276 & $-$55,1 & $-$69,1 & 68,8 & $-$6,07 \\
    4 & $\alpha_v = 0^\circ$, t1, dir. & $-0{,}3$ & 3,78 & 0,179 & 134 & 8,35 & 127 & 34,3 & 10,7 & 29,2 & 2,57 \\
    5 & $\alpha_v = 0^\circ$, t2, esq. & $+0{,}8$ & $-$2,70 & $-$1,98 & $-$93,7 & 8,76 & $-$92,4 & $-$27,9 & $-$21,8 & $-$25,4 & $-$2,24 \\
    6 & $\alpha_v = 0^\circ$, t2, esq. & $-0{,}3$ & 5,21 & 5,93 & 311 & 10,4 & 311 & 61,4 & 58,0 & $-$57,8 & 5,10 \\
    7 & $\alpha_v = 0^\circ$, t2, dir. & $+0{,}8$ & $-$1,98 & $-$2,70 & $-$93,9 & 11,5 & $-$92,6 & $-$23,4 & $-$25,5 & 25,4 & $-$2,24 \\
    8 & $\alpha_v = 0^\circ$, t2, dir. & $-0{,}3$ & 5,93 & 5,21 & 311 & 9,92 & 310 & 65,9 & 54,4 & 58,3 & 5,14 \\
    9 & $\alpha_v = 90^\circ$, b1 & $+0{,}8$ & $-$5,57 & $-$5,57 & $-$257 & 10,2 & $-$257 & $-$62,2 & $-$56,3 & $-$56,4 & $-$4,97 \\
    10 & $\alpha_v = 90^\circ$, b1 & $-0{,}3$ & 2,34 & 2,34 & 146 & 10,2 & 146 & 27,1 & 23,5 & 23,7 & 2,09 \\
    11 & $\alpha_v = 90^\circ$, b2 & $+0{,}8$ & $-$4,14 & $-$4,14 & $-$184 & 10,2 & $-$184 & $-$46,0 & $-$41,8 & $-$41,8 & $-$3,69 \\
    12 & $\alpha_v = 90^\circ$, b2 & $-0{,}3$ & 3,78 & 3,78 & 219 & 10,2 & 219 & 43,4 & 38,1 & 38,2 & 3,37 \\
    13 & $\alpha_v = 90^\circ$, b3 & $+0{,}8$ & $-$1,26 & $-$1,26 & $-$37,8 & 8,00 & $-$37,5 & $-$13,5 & $-$12,8 & 12,7 & $-$1,12 \\
    14 & $\alpha_v = 90^\circ$, b3 & $-0{,}3$ & 6,65 & 6,65 & 365 & 10,2 & 365 & 75,9 & 67,1 & 67,4 & 5,94 \\
\midrule
    \multicolumn{12}{@{}l}{\textit{Vento de sucção, $0{,}9\,G + 1{,}5\,W$}} \\
    \textbf{15} & \textbf{$\alpha_v = 0^\circ$, t1, esq.} & \textbf{$+0{,}8$} & \textbf{$-$9,71} & \textbf{$-$6,11} & \textbf{$-$389} & \textbf{9,11} & \textbf{$-$385} & \textbf{$-$100} & \textbf{$-$70,7} & \textbf{$-$89,1} & \textbf{$-$7,86} \\
    16 & $\alpha_v = 0^\circ$, t1, esq. & $-0{,}3$ & $-$1,80 & 1,80 & 34,1 & 13,8 & 18,3 & $-$10,7 & 9,12 & $-$9,10 & $-$0,802 \\
    17 & $\alpha_v = 0^\circ$, t1, dir. & $+0{,}8$ & $-$6,11 & $-$9,71 & $-$391 & 11,2 & $-$386 & $-$77,7 & $-$89,1 & 88,7 & $-$7,82 \\
    18 & $\alpha_v = 0^\circ$, t1, dir. & $-0{,}3$ & 1,80 & $-$1,80 & 32,7 & 6,48 & 17,1 & 11,7 & $-$9,20 & 9,17 & $-$0,808 \\
    19 & $\alpha_v = 0^\circ$, t2, esq. & $+0{,}8$ & $-$4,68 & $-$3,96 & $-$202 & 9,67 & $-$202 & $-$50,5 & $-$41,8 & $-$45,5 & $-$4,01 \\
    20 & $\alpha_v = 0^\circ$, t2, esq. & $-0{,}3$ & 3,24 & 3,96 & 201 & 10,5 & 201 & 38,9 & 38,1 & $-$37,9 & 3,34 \\
    21 & $\alpha_v = 0^\circ$, t2, dir. & $+0{,}8$ & $-$3,96 & $-$4,68 & $-$203 & 10,6 & $-$202 & $-$46,0 & $-$45,4 & 45,3 & $-$3,99 \\
    22 & $\alpha_v = 0^\circ$, t2, dir. & $-0{,}3$ & 3,96 & 3,24 & 201 & 9,77 & 201 & 43,3 & 34,4 & 38,2 & 3,37 \\
    23 & $\alpha_v = 90^\circ$, b1 & $+0{,}8$ & $-$7,55 & $-$7,55 & $-$367 & 10,2 & $-$367 & $-$84,8 & $-$76,3 & $-$76,4 & $-$6,73 \\
    24 & $\alpha_v = 90^\circ$, b1 & $-0{,}3$ & 0,359 & 0,359 & 36,0 & 10,2 & 36,0 & 4,55 & 3,59 & 3,67 & 0,323 \\
    25 & $\alpha_v = 90^\circ$, b2 & $+0{,}8$ & $-$6,11 & $-$6,11 & $-$294 & 10,2 & $-$294 & $-$68,5 & $-$61,7 & $-$61,8 & $-$5,45 \\
    26 & $\alpha_v = 90^\circ$, b2 & $-0{,}3$ & 1,80 & 1,80 & 109 & 10,2 & 109 & 20,8 & 18,1 & 18,2 & 1,61 \\
    27 & $\alpha_v = 90^\circ$, b3 & $+0{,}8$ & $-$3,24 & $-$3,24 & $-$147 & 10,2 & $-$147 & $-$36,1 & $-$32,7 & $-$32,7 & $-$2,88 \\
    28 & $\alpha_v = 90^\circ$, b3 & $-0{,}3$ & 4,67 & 4,67 & 256 & 10,2 & 256 & 53,3 & 47,1 & 47,3 & 4,17 \\
\midrule
    \multicolumn{12}{@{}l}{\textit{Sobrecarga dominante, $1{,}35\,G + 1{,}5\,Q$}} \\
    \textbf{29} & \textbf{sem vento} &  & \textbf{9,38} & \textbf{9,38} & \textbf{505} & \textbf{10,2} & \textbf{505} & \textbf{107} & \textbf{94,7} & \textbf{95,0} & \textbf{8,37} \\
\bottomrule
\end{tabular}}{Elaborado pelo autor (2026)}{an-esforcos}{}{}

Na tabela, $q_{esq}$ e $q_{dir}$ são as cargas de cálculo por água, positivas no sentido da gravidade; $x$ é a posição do momento máximo $M_{max}$ a partir do apoio esquerdo; $M_{ap,d}$ é o momento na cumeeira; $R_A$ e $R_B$ são as reações verticais nos apoios; $V_{max}$ e $N_{max}$ são o cortante e o esforço normal no eixo da peça, cujos extremos ocorrem nos apoios. As linhas em negrito são as combinações governantes. As combinações de vento têm duração curta e a de sobrecarga, duração média, o que define o $k_{mod}$ de cada verificação.
